%%% XeLaTeX-article %%%
%# -*- coding: utf-8 -*-
%!TEX encoding = UTF-8 Unicode
%!TEX TS-program = xelatex  
%---------------------虽然加了%还是要保留！

\documentclass[17pt]{beamer}
\mode<presentation>
{
\usetheme[width=40pt]{Hannover}
\usecolortheme[]{dove}
\usefonttheme[]{structurebold}
\setbeameroption{hide notes}
}

\usepackage{fontspec}
\usepackage{polyglossia}
\setmainfont{Arial} %设置主字体
\newfontfamily\sanskritfont[Script=Devanagari,Mapping=romantodevanagari,Scale=1.15]{Sanskrit 2003}             %输出天城体
%\newfontfamily\sanskritfont[Mapping=tex-text]{Times New Roman}              %输出转写
\doublehyphendemerits=-10000
\newcommand{\skt}[1]{{\sanskritfont{#1}}} %输出天城体
\newcommand{\skttrans}[1]{{\skt{#1}~#1}}  %输出天城体和转写
%----------------------------------------------------设置梵文输入方法 danda । ॥

\usepackage[UTF8,fontset=windows]{ctex}
\usepackage{amsmath}
%----------------------------------------------------设置中文环境

\usepackage{graphicx}
\usepackage{flafter} 
\graphicspath{{pic/}}
\usepackage{booktabs} 
\usepackage{nicematrix}
\newenvironment{indentlist}
  {\begin{list}{}{\setlength{\leftmargin}{2em}\setlength{\itemsep}{0.5em}}}
  {\end{list}}
%-----------------------------------------插图表格

\usepackage{hyperref} 
\usepackage[dvipsnames]{xcolor}
\usepackage{colortbl}
\definecolor{light-gray}{gray}{0.9}
%------------------------------颜色

\newcommand{\verbroot}[1]{\textcolor{red}{$\sqrt{}$#1}}
\newcommand{\sktroot}[1]{{\verbroot{\skt{#1}}}}
\newcommand{\skttransroot}[1]{{\sktroot{#1}~\textcolor{red}{#1}}}

\newcommand{\nounstem}[1]{\textcolor{red}{#1\nobreakdash-}}
\newcommand{\sktnounstem}[1]{{\textcolor{red}{\skt{#1\nobreakdash-}}}}
\newcommand{\skttransnounstem}[1]{{\sktnounstem{#1}~\nounstem{#1}}}

\newcommand{\verbstem}[1]{\textcolor{blue}{#1\nobreakdash-}}
\newcommand{\sktverbstem}[1]{{\textcolor{blue}{\skt{#1\nobreakdash-}}}}
\newcommand{\skttransverbstem}[1]{{\sktverbstem{#1}~\verbstem{#1}}}

\newcommand{\wordending}[1]{\textcolor{Orange}{\nobreakdash-#1}}
\newcommand{\sktending}[1]{{\textcolor{Orange}{\skt{-#1}}}}
\newcommand{\skttransending}[1]{{\sktending{#1}~\wordending{#1}}}

\newcommand{\fullpada}[1]{\textcolor{OliveGreen}{#1}}
\newcommand{\sktpada}[1]{{\textcolor{OliveGreen}{\skt{#1}}}}
\newcommand{\skttranspada}[1]{{\sktpada{#1}~\fullpada{#1}}}

\newcommand{\pratyaya}[1]{\textcolor{Plum}{#1}}
\newcommand{\sktpratyaya}[1]{{\textcolor{Plum}{\skt{#1}}}}
\newcommand{\skttranspratyaya}[1]{{\sktpratyaya{#1}~\pratyaya{#1}}}

\newcommand{\reconstruction}[1]{\textcolor{gray}{*#1}}
\newcommand{\fullsentence}[1]{\textcolor{MidnightBlue}{#1}}
\newcommand{\sktsentence}[1]{\textcolor{MidnightBlue}{\skt{#1}}}

\newcommand{\veryimportant}[1]{\textcolor{red}{#1}}
\newcommand{\important}[1]{\textcolor{blue}{#1}}
\newcommand{\notsoimportant}[1]{\textcolor{gray}{#1}}
\newcommand{\dicturl}[2]{\href{#1}{\mbox{\texttt{#2}}}}
%-------------------------------------------词根等标颜色

\title{{梵语提高}}
\subtitle{1. 梵语工具}
\author[张雪杉]{文学院~~张雪杉 \\ zhangxueshan@sdnu.edu.cn}
\date{}

\begin{document}	

\begin{frame}
  \titlepage
\end{frame}

\begin{frame}
  \frametitle{本节内容}
  \small
  \tableofcontents
\end{frame}

\section{西文资源}

\begin{frame}{\insertsection }
  \begin{itemize}
    \item 词典：Cologne Sanskrit Lexicon
    \item 变形：Sanskrit Heritage
  \end{itemize}
\end{frame}

\subsection{Cologne 科隆}
\begin{frame}{Cologne Sanskrit Lexicon}
  \small
  \begin{itemize}
    \item 网址：\dicturl{https://www.sanskrit-lexicon.uni-koeln.de/}{sanskrit-lexicon.uni-koeln.de}
    \item 性质：学术检索入口，聚合多部历史词典并支持互跳。
    \item 底本：MW，PW，PWG，Apte。
  \end{itemize}
\end{frame}

\begin{frame}{MW. Monier-Williams}
  \small
  \begin{itemize}
    \item Monier\mbox{-}Williams Sanskrit\mbox{-}English Dictionary
    \item 版本：1899 年增补版，多次重印。
    \item 特点：义项系统完整，注释其他印欧语言同源词，最常用。
  \end{itemize}
\end{frame}

\begin{frame}{MW. 缩略语}
  \small
  \begin{itemize}
    \item List of Works and Authors. pp. xxxiii-xxxiv.\\
    \begin{itemize}
      \item Ṛ(ig-)V(eda, refered to as RV).
      \item Kāthās(aritsāgara).
    \end{itemize}
    \item Symbols and Abbreviations pp. xxxiv-xxxv.\\
    \begin{itemize}
      \item Ā. = Ātmane-pada.
      \item abl. = ablative case.
    \end{itemize}    
  \end{itemize}
\end{frame}

\begin{frame}{PWG./PW. \\Böhtlingk德语词典}
  \small
  \begin{itemize}
    \item Böhtlingk and Roth Grosses Petersburger Wörterbuch (PWG, 1855)
    \item Böhtlingk Sanskrit-Wörterbuch in kürzerer Fassung (PW, 1879)
    \item 特点：历史分层信息丰富，适合语义史与文献义项比对。
  \end{itemize}
\end{frame}

\begin{frame}{AP. Apte 词典}
  \small
  \begin{itemize}
    \item Practical Sanskrit-English Dictionary
    \item 版本：1957修订增补本
    \item 用途：条理清楚，教学友好，适合与 MW/PWG 交叉核对。
  \end{itemize}
\end{frame}

\subsection{Sanskrit Heritage}
\begin{frame}{Dictionnaire Héritage du Sanskrit}
  \small
  \begin{itemize}
    \item 网址：\dicturl{https://sanskrit.inria.fr/}{sanskrit.inria.fr}
    \item 性质：梵语形态分析与分词工具。
    \item 特点：从屈折形式反查词头，并提供形态解析路径。
  \end{itemize}
\end{frame}

\section{中文资源}

\begin{frame}{\insertsection }
  \begin{itemize}
    \item 综合：佛研资讯
    \item 顺手：梵佛研
    \item 也常用：梵和大辞典
  \end{itemize}
\end{frame}

\subsection{佛研资讯}
\begin{frame}{佛研资讯}
  \small
  \begin{itemize}
    \item 网址：\dicturl{https://foyan.online/}{foyan.online}
    \item 性质：佛教数字资源聚合入口。
    \item 内容结构：包含数字佛典平台（CBETA、SAT）与电子词典平台（DDB）等。
    \item 特点：新且较详细，源自2024综述文章。
  \end{itemize}
\end{frame}

\begin{frame}{CBETA 数字佛典平台}
  \small
  \begin{itemize}
    \item 网址：\dicturl{https://cbetaonline.dila.edu.tw/}{cbetaonline.dila.edu.tw}
    \item 性质：汉文佛典数据库与检索平台。
    \item 用途：定位术语在不同译本、不同语境中的用例与译法。
  \end{itemize}
\end{frame}

\begin{frame}{SAT 数字佛典平台}
  \small
  \begin{itemize}
    \item 网址：\dicturl{https://21dzk.l.u-tokyo.ac.jp/SAT/}{21dzk.l.u-tokyo.ac.jp/SAT}
    \item 性质：大藏经检索与并行对读平台。
    \item 用途：精确检索汉译佛典术语，补充 CBETA 的检索路径。
  \end{itemize}
\end{frame}

\begin{frame}{DDB 电子佛教词典}
  \small
  \begin{itemize}
    \item 网址：\dicturl{https://www.buddhism-dict.net/ddb/}{buddhism-dict.net/ddb}
    \item 性质：佛教术语英文释义、汉字词条与跨语言对照。
    \item 特点：主要是以汉查梵？
  \end{itemize}
\end{frame}


\subsection{梵佛研}
\begin{frame}{梵佛研}
  \small
  \begin{itemize}
    \item 网址：\dicturl{https://fanfoyan.com/}{fanfoyan.com}
    \item 资源：梵语宝典和梵佛词典。
    \item 特点：主要是Monier-Williams，部分词有变格变位表和中文义项。
  \end{itemize}
\end{frame}

\subsection{荻原梵和大辞典}
\begin{frame}{漢譯對照梵和大辭典}
  \small
  \begin{itemize}
    \item 网址：\dicturl{https://www.manduuka.net/sanskrit/dic/ogdic.htm}{manduuka.net/sanskrit/dic/ogdic.htm}
    \item 编者：荻原雲來
    \item 性质：梵日漢對照佛教辭典。
    \item 特点：只有上册可查。
  \end{itemize}
\end{frame}

\section{下节预告}

\begin{frame}{\insertsection }
  \begin{itemize}
    \item 上学期作业：第32章练习3
    \item 第33章：愿望动词和必要分词
  \end{itemize}
\end{frame}  

\end{document}
